\documentclass[10pt,a4paper]{amsart}
\usepackage[margin=19mm]{geometry}
\usepackage{amsmath,amssymb,mathtools}
\usepackage[scr=rsfs,cal=euler]{mathalfa}
\usepackage{xfrac}
\usepackage[unicode,hidelinks,bookmarks=false]{hyperref}
\newcommand{\ie}{i.e.\ }
\newcommand{\cf}{cf.\ }
\newcommand{\resp}{respectively\ }
\newcommand{\Subsection}{\subsection}
\providecommand{\nobreakdash}{\nobreak\hbox{-}}
\setlength{\emergencystretch}{2em}
\multlinegap0pt
\usepackage{bm}
\usepackage{bbm}
\usepackage{dsfont}
\usepackage{xspace}
\usepackage{cancel}
\usepackage{mathtools}
\usepackage{amsthm}
\makeatletter
\@ifundefined{theorem}{\newtheorem{theorem}{Theorem}}{}
\makeatother
\title[Permutation-based Strategies for Labeled Chip-Firing on $k$-]{Permutation-based Strategies for Labeled Chip-Firing on $k$-ary Trees}
\author{Ryota Inagaki, Tanya Khovanova, Austin Luo}
\date{Source: arXiv:2503.09577v4 (CC BY 4.0)}
\begin{document}
\maketitle

\noindent\emph{Source and licence.} Permutation-based Strategies for Labeled Chip-Firing on $k$-ary Trees. Version arXiv:2503.09577v4 (CC BY 4.0), distributed under the Creative Commons Attribution 4.0 International licence, as recorded in the arXiv licence metadata for this exact version and shown on the abstract page https://arxiv.org/abs/2503.09577v4. Full rights notice: licenses/arxiv-2503.09577-CC-BY-4.0.txt.\par\smallskip

\noindent\textit{Editorial scope.} Counted unit: the Theorem whose source label is \texttt{thm:Lehmer} in the article (source0.tex lines 326--329, source label \texttt{thm:Lehmer}), delivered together with the article's own complete original proof of it (source0.tex lines 331--365, one \texttt{proof} environment of 35 lines). No mathematics is written, repaired or re-derived: statement and proof are the article's own text line for line. Required context retained uncounted: none. Every definition, display and label of those lines is retained unchanged. Omitted from this package and not redistributed: 1--325, 330--330, 366--654. Nothing inside a retained excerpt is omitted or altered: every retained body is delivered exactly as the article's lines are written, so no omission receipt of any kind accompanies this package. Citations: the retained excerpts carry no citation command at all, so no bibliography entry of the article is transcribed and no reference is redistributed. Reference adaptation: none was needed and none was made; every reference inside the delivered unit resolves inside it, so no unresolved reference or citation is produced. Printed slips kept literal: no printed word, symbol, hypothesis or number of the counted statement or of its proof was altered. Layout and class: the article's own class is \texttt{dmtcs-episciences} with packages including inputenc, subfigure, natbib, amsmath, amssymb, mathalpha, amsthm; the wrapper is the lane's pinned \texttt{amsart} editorial wrapper at 10pt on A4, which restates the theorem-like environments the retained text needs and adds the article's own macro definitions that the retained text needs (none), with extra wrapper packages (bm, bbm, dsfont, xspace, cancel, mathtools). The delivered numbering is the unit's own. Assets: no retained span contains a figure environment, a graphics-inclusion command, a PDF-page inclusion, a table or a TikZ picture; no image file is copied into this package, the package declares no asset dependency and no image file is redistributed with it. Licence: version arXiv:2503.09577v4 of this article is distributed under the Creative Commons Attribution 4.0 International licence, as recorded in the arXiv licence metadata for this exact version and shown on the abstract page https://arxiv.org/abs/2503.09577v4. A full rights notice is delivered at licenses/arxiv-2503.09577-CC-BY-4.0.txt. Characters: every retained excerpt is pure ASCII, so the delivered engine, XeLaTeX with its default Latin Modern text font, reports no missing character.
\begin{theorem}\label{thm:Lehmer}
Let $k \geq 2$, $n \in \mathbb{N}$, and $w$ be any permutation in $S_n$. Let $c_w$ be the Lehmer code of permutation $w$. Then the number of inversions in $\mathcal{C}_{k, n, w}$ is
\[I(k, n, w) = \binom{k}{2} \sum_{i=1}^{n}\binom{k^{(c_w)_{i}}}{2}k^{2n-i-2(c_w)_i-1} = \frac{k^n(k-1)}{4}\sum_{i=1}^n(k^{n-i} -k^{n-i-(c_w)_i}).\]
\end{theorem}

\begin{proof}
Let $w = w_1w_2\dots w_n$. First, we count the inversions in $\mathcal{C}_{k, n, w}$ that result from two chips that end up in subtrees rooted at different children of the root vertex. Consider that when sorting chips by the $w_1$th most significant digit in the $k$-ary representation, we know that the $j$th child from the left contains the set of chips
$$S_j = \left\{k^{n-w_1}(j-1) + \sum_{m \in \{0,1, \dots, n-1\}\setminus \{n-w_1\}}d_m k^m: d_m \in \{0,1, \dots, k-1\} \right\}.$$
Therefore, for each $j', j \in [k]$ such that $j' > j$, we find that for each chip $x = k^{n-w_1}(j-1)+\sum_{m \in \{0, 1, \dots, n-1\}\setminus \{n-w_1\}}d_mk^m$ in the $j$th child of the root, the set of chips in the $j'$th child of the root that are less than $x$ is
$$\left\{k^{n-w_1}(j-1) + \sum_{m \in \{0,1, \dots, n-1\}\setminus \{n-a_1\} }d_m'k^m: \sum_{m = n-w_1+1}^{n-1}d_m'k^{m} < \sum_{m = n-w_1+1}^{n-1} d_mk^{m} \right\}.$$ The cardinality of this set is $k^{n-w_1}\sum_{m=0}^{w_1-2}d_{m+n-w_1+1}k^m $. Therefore, the number of inversions in $C_{k, n, w}$ that result from a chip sent to the $j'$th child of the root and another sent to the $j$th child of the root is
$$\binom{k^{c_i}}{2}k^{2n-2w_1}.$$
Therefore, the total number of inversions that result from two chips ending up in the subtrees rooted at different children of the root is
$$\sum_{j=1}^{k-1}(k-i)\binom{k^{w_1-1}}{2} (k^{n-w_1})^2.$$
Since $(c_w)_1 = w_1-1$ by definition of Lehmer code of $w$, we find that this quantity is equal to
$$\sum_{j=1}^{k-1}(k-i)\binom{k^{(c_w)_1} }{2} (k^{n-(c_w)_1-1})^2.$$

Note that after the root cannot fire, we find that there are $k^{n-1}$ chips on each root's child. When ignoring the $w_1$th most significant digit, taking relative order $w'$ of $w_2w_3\dots w_n$, the firing of the $k^{n-1}$ chips on each child of the root will yield the same number of inversions as firing $k^{n-1}$ using strategy $w'$. The total number of inversions obtained from two chips that are in the same subtree of the root is $I(k, n-1, w')$. Thus,
\begin{equation}\label{eq:RecursiveRelation}
I(k, n, w) = \sum_{i=1}^{k-1}(k-i)\binom{k^{(c_w)_1} }{ 2} (k^{n-(c_w)_1-1})^2 + k I(k, n-1, w').
\end{equation}

We now prove that for fixed $k \geq 2$, for all $n \in \mathbb{N}$ and $w \in S_n$, we have
\[I(k, n, w) = \binom{k}{2} \sum_{i=1}^{n}\binom{k^{(c_w)_{i}} }{ 2}k^{2n-i-2(c_w)_i-1}\]
via induction on $n$. For our base case, $n=0$, we know that $\mathcal{C}_{k, 0, w}$ only consists of one chip. Hence $I(k,0, w) = 0$. Now consider the inductive step. 

We shall show that for any $k \geq 2$ and $w \in S_{n_0+1}$, we have
\[I(k, n_0+1, w) = \binom{k}{2} \sum_{i=1}^{n_0+1} \binom{k^{(c_w)_i}}{2}k^{2n_0+2-i-2(c_w)_i-1}\]
assuming that for any $w' \in S_{n_0}$ we have $$I(k, n_0, w') = \binom{k}{2} \sum_{i=1}^{n_0}\binom{k^{(c_{w'})_{i}} }{ 2}k^{2n-i-2(c_{w'})_i-1}.$$
By Eq.~(\ref{eq:RecursiveRelation}), we have $I(k, n+1, w) = \binom{k}{2}\binom{k^{(c_w)_1}}{2}k^{2n+2-c_1-1} + k I(k, n_0, w')$, where $w' \in S_{n_0} $ is the result of taking the relative order of $w_2w_3\dots w_{n_0+1}$. Since $(c_{w'})_{j} = (c_{w})_{j+1}$ for each $j \in [n_0]$, the inductive hypothesis tells us that \begin{equation*}
    \begin{split}
        I(k, n+1, w) &= \binom{k}{2}\binom{k^{(c_w)_1}}{2}k^{2n+2-(c_w)_1-1} + k I(k, n_0, w') \\ &= \binom{k}{2}\binom{k^{c_1}}{2}k^{2n+2-c_1-1} +k \binom{k}{2}\sum_{i=1}^{n_0} \binom{k^{(c_{w})_{i+1}}}{2} k^{2n_0 - i-2(c_w)_{i+1}-1} \\ &= \binom{k}{2}\sum_{i=1}^{n_0+1}\binom{k^{(c_w)_i}}{2}k^{2n_0+1-i-2(c_w)_{i}}.
    \end{split}
\end{equation*} This proves the inductive hypothesis and hence completes the proof of the first formula.
For the second formula, we have
\begin{equation*}
\begin{split}
I(k, n, w) &= \binom{k}{2}\sum_{i=1}^n\binom{k^{(c_w)_i}}{2}k^{2n-i-2(c_w)_i-1} = \frac{k(k-1)}{4}\sum_{i=1}^n (k^{2(c_w)_i}-k^{(c_w)_i})k^{2n-i-2(c_w)_i-1} \\
&= \frac{k-1}{4}\sum_{i=1}^n (k^{2n-i}-k^{2n-i-(c_w)_i}) = \frac{k^n(k-1)}{4}\sum_{i=1}^n(k^{n-i} -k^{n-i-(c_w)_i}).
\end{split}
\end{equation*}\end{proof}

\end{document}
